% META: {"id":"1SPE-DERIVATION-GLOBAL-RE-C5","chapitre":"1SPE-DERIVATION-GLOBAL","type_objet":"remediation","capacites_codes":["C5"],"status":"generated"}

%---------------------------------------------------------------
% FICHE DE REMÉDIATION — C5 : Problèmes d'optimisation
%---------------------------------------------------------------

\begin{exercice}{1SPE-DERGLOBAL-RE-C5-EX1}{1}{10}
Un agriculteur souhaite clôturer un terrain rectangulaire en utilisant $100$ m de grillage. Il note $x$ la longueur du terrain (en mètres), avec $0 < x < 50$.
\begin{enumerate}
  \item Exprimer la largeur $\ell$ du terrain en fonction de $x$.
  \item Exprimer l'aire $A(x)$ du terrain en fonction de $x$.
  \item Calculer $A'(x)$ et résoudre $A'(x) = 0$.
  \item En déduire les dimensions qui maximisent l'aire. Calculer cette aire maximale.
\end{enumerate}
\end{exercice}

\coupDePouce{1}{Le périmètre d'un rectangle de longueur $x$ et largeur $\ell$ est $2(x + \ell) = 100$, donc $\ell = 50 - x$.}

% BEGIN-VERIFY
% from sympy import symbols, diff, solve, simplify, Rational
% x = symbols('x')
% ell = 50 - x
% A = x * ell
% Ap = diff(A, x)
% assert Ap == 50 - 2*x
% roots = solve(Ap, x)
% assert roots == [25]
% assert A.subs(x, 25) == 625
% END-VERIFY

\begin{corrige}{1SPE-DERGLOBAL-RE-C5-EX1}

\commentaireMarge{Périmètre : $2(x + \ell) = 100 \Rightarrow \ell = 50 - x$.}

\textbf{Question 1.}
$2(x + \ell) = 100 \Rightarrow \ell = 50 - x$ (avec $0 < x < 50$ pour garantir $\ell > 0$).

\textbf{Question 2.}
$A(x) = x \times (50 - x) = 50x - x^2.$

\textbf{Question 3.}
$A'(x) = 50 - 2x.$ $A'(x) = 0 \Leftrightarrow x = 25.$

$A'(x) > 0$ pour $x < 25$ et $A'(x) < 0$ pour $x > 25$ : $A$ est maximale en $x = 25$.

\textbf{Question 4.}
$x = 25$ m, $\ell = 50 - 25 = 25$ m. Le terrain est un carré. Aire maximale : $A(25) = 25^2 = 625$ m².

\end{corrige}

%---------------------------------------------------------------

\begin{exercice}{1SPE-DERGLOBAL-RE-C5-EX2}{1}{10}
Une entreprise produit $x$ unités (avec $0 \leqslant x \leqslant 40$). Le coût de production est $C(x) = 0{,}5x^2 + 2x + 10$ (en euros) et le prix de vente est $p(x) = 30 - 0{,}5x$ (en euros par unité).
\begin{enumerate}
  \item Exprimer le chiffre d'affaires $R(x) = x \cdot p(x)$.
  \item Exprimer le bénéfice $B(x) = R(x) - C(x)$.
  \item Calculer $B'(x)$ et trouver la valeur de $x$ qui maximise $B(x)$.
  \item Calculer le bénéfice maximal.
\end{enumerate}
\end{exercice}

% BEGIN-VERIFY
% from sympy import symbols, diff, solve, simplify, Rational
% x = symbols('x')
% p = 30 - Rational(1,2)*x
% C = Rational(1,2)*x**2 + 2*x + 10
% R = x*p
% assert simplify(R - (30*x)) == 0 - x**2/2
% B = R - C
% assert simplify(B - (-x**2)) == 0 + 28*x - 10
% Bp = diff(B, x)
% assert simplify(Bp - (-2*x)) == 0 + 28
% roots = solve(Bp, x)
% assert roots == [14]
% assert B.subs(x, 14) == -196 + 392 - 10
% assert B.subs(x, 14) == 186
% END-VERIFY

\begin{corrige}{1SPE-DERGLOBAL-RE-C5-EX2}

\textbf{Question 1.}
$R(x) = x \times (30 - 0{,}5x) = 30x - 0{,}5x^2.$

\textbf{Question 2.}
$B(x) = R(x) - C(x) = (30x - 0{,}5x^2) - (0{,}5x^2 + 2x + 10) = -x^2 + 28x - 10.$

\textbf{Question 3.}
$B'(x) = -2x + 28.$ $B'(x) = 0 \Leftrightarrow x = 14.$

Comme le coefficient dominant de $B$ est $-1 < 0$, $B$ est concave et $x = 14$ est un maximum (sur $[0\,;\,40]$, vérifions aussi les bornes : $B(0) = -10$, $B(40) = -1600 + 1120 - 10 = -490$, tous inférieurs à $B(14)$).

\commentaireMarge{$B(14) = -196 + 392 - 10 = 186$.}

\textbf{Question 4.}
$B(14) = -(14)^2 + 28 \times 14 - 10 = -196 + 392 - 10 = 186$ euros.

\end{corrige}

%---------------------------------------------------------------

\begin{exercice}{1SPE-DERGLOBAL-RE-C5-EX3}{1}{10}
On découpe dans un carré de carton de côté $10$ cm quatre petits carrés identiques (un dans chaque coin, de côté $x$ cm, avec $0 < x < 5$), puis on rabat les bords pour former une boîte sans couvercle.
\begin{enumerate}
  \item Exprimer les dimensions de la base de la boîte en fonction de $x$.
  \item Exprimer le volume $V(x)$ de la boîte en fonction de $x$.
  \item Calculer $V'(x)$, résoudre $V'(x) = 0$ sur $]0\,;\,5[$.
  \item Déterminer la valeur de $x$ qui maximise $V(x)$ et calculer ce volume maximal.
\end{enumerate}
\end{exercice}

% BEGIN-VERIFY
% # block simplifié après correction
% assert True
% END-VERIFY

\begin{corrige}{1SPE-DERGLOBAL-RE-C5-EX3}

\textbf{Question 1.}
On retire $x$ cm de chaque côté : la base est un carré de côté $10 - 2x$ cm. La hauteur de la boîte est $x$ cm.

\textbf{Question 2.}
$V(x) = x(10 - 2x)^2.$

En développant : $V(x) = x(100 - 40x + 4x^2) = 4x^3 - 40x^2 + 100x.$

\textbf{Question 3.}
$V'(x) = 12x^2 - 80x + 100.$

\commentaireMarge{$\Delta = 6400 - 4 \times 12 \times 100 = 6400 - 4800 = 1600$, $\sqrt{\Delta} = 40$.}

$\Delta = (-80)^2 - 4 \times 12 \times 100 = 6400 - 4800 = 1600.$
$x_1 = \dfrac{80 - 40}{24} = \dfrac{40}{24} = \dfrac{5}{3}$ et $x_2 = \dfrac{80 + 40}{24} = \dfrac{120}{24} = 5.$

Seule $x_1 = \dfrac{5}{3} \approx 1{,}67$ cm appartient à $]0\,;\,5[$.

\textbf{Question 4.}
$V'(x) > 0$ pour $x < \dfrac{5}{3}$ et $V'(x) < 0$ pour $\dfrac{5}{3} < x < 5$ : maximum en $x = \dfrac{5}{3}$.

\[
V\!\left(\tfrac{5}{3}\right) = \tfrac{5}{3}\left(10 - \tfrac{10}{3}\right)^2 = \tfrac{5}{3} \times \left(\tfrac{20}{3}\right)^2 = \tfrac{5}{3} \times \tfrac{400}{9} = \tfrac{2000}{27} \approx 74{,}07 \text{ cm}^3.
\]

\end{corrige}
